% ============================================================
%  ch4_III  第三章 无水平结构的紧化；第四章 紧化上的截断与水平结构（开头）
%  原书 pp. 19--27，即 PDF 第 20--28 页
% ============================================================

\section{无水平结构的紧化}

$X$ 是 $\mathbb{F}_q$ 上的光滑（几何连通）射影曲线，$r\ge 2$ 为整数。

\par\medskip\noindent\textbf{所需如下.}\ 
\begin{itemize}
\item 已经定义了开集
\[
\mathrm{Cht}^{r,d,\overline{d}\le p}\subset\mathrm{Cht}^{r,d}\subset\mathrm{Cht}^{r}
\]
它们在 $X\times X$ 上有限型且光滑。
\item 要构造叠 $\overline{\mathrm{Cht}}{}^{r,d,\overline{d}\le p}$，使得
\item \origpage{20}$\overline{\mathrm{Cht}}{}^{r,d,\overline{d}\le p}$ 以 $\mathrm{Cht}^{r,d,\overline{d}\le p}$ 为稠密开集，
\item 结构态射 $\mathrm{Cht}^{r,d,\overline{d}\le p}\to X\times X$ 延拓为紧合光滑态射
\[
\overline{\mathrm{Cht}}{}^{r,d,\overline{d}\le p}\to X\times X,
\]
\item 边界 $\overline{\mathrm{Cht}}{}^{r,d,\overline{d}\le p}-\mathrm{Cht}^{r,d,\overline{d}\le p}$ 是相对正规交叉除子，
\item 边界各分层在 $X\times X$ 上的 $\ell$-进上同调是“$r$-可忽略”的，
\item 在 $\overline{\mathrm{Cht}}{}^{r,\overline{d}\le p}/a^{\mathbb{Z}}\cong\coprod_{1\le d\le r|\deg(a)|}\overline{\mathrm{Cht}}{}^{r,d,\overline{d}\le p}$ 中，Hecke 对应在“其不动点附近”保持开集 $\mathrm{Cht}^{r,\overline{d}\le p}/a^{\mathbb{Z}}$ 不变。
\end{itemize}

\par\medskip\noindent\textbf{回顾.}\ 设 $S$ 是 $\mathbb{F}_q$ 上的概形。$S$ 上秩 $r$ 的 chtouca 由下列各项构成：

(i) $X\times S$ 上秩 $r$ 的局部自由丛 $\mathcal{E}$，

(ii) $\mathcal{E}$ 的一个修正 $\mathcal{E}\hookrightarrow\mathcal{E}'\hookleftarrow\mathcal{E}''$，

(iii) 一个同构 ${}^\tau\mathcal{E}=(\mathrm{Id}_X\times\mathrm{Frob}_S)^*\mathcal{E}\xrightarrow{\sim}\mathcal{E}''$。

\par\medskip\noindent\textbf{想法.}\ 放宽模问题：把 (iii) 中的“同构”换成“完整同态”。

\par\medskip\noindent\textbf{完整同态.}\ 先设：$A$ 为离散赋值环，$\pi_A$ 为一致化子，
$K_A=\operatorname{Frac}(A)\qquad \kappa_A=A/(\pi_A)$。
设 $v\in\mathrm{GL}_r(K_A)\cap(M_r-\{0\})(A)$。
我们要理解 $u$\footnote{译注：此处原文作 $u$，而上文所设的矩阵记为 $v$，原文如此。[译者按（上下文推断，无外部依据）：下文同此。]} 在 $\kappa_A$ 之上如何退化。按初等因子理论，可以写成
\[
v=g'\begin{pmatrix}
1 & & & & \\
& \pi_A^{d_1} & & & \\
& & \pi_A^{d_1+d_2} & & 0\\
& 0 & & \ddots & \\
& & & & \pi_A^{d_1+\cdots+d_{r-1}}
\end{pmatrix}g
\]

\origpage{21}其中 $d_1,d_2,\ldots,d_{r-1}\in\mathbb{N}$，$g,g'\in\mathrm{GL}_r(A)$。
设 $\underline{r}$ 是唯一的划分 $(r=r_1+\cdots+r_k)$，使得对 $1\le s<r$ 有
\[
d_s=0\iff s\in\underbrace{\{0,r_1,\ldots,r_1+\cdots+r_{k-1}\}}_{\underline{r}^-}\cap\underbrace{\{r_1,r_1+r_2,\ldots,r\}}_{\underline{r}^+}.
\]
对每个 $s$，$1\le s\le r$，同态
\[
\Bigl(\prod_{t<s}\pi_A^{-d_t(s-t)}\Bigr)\cdot\Lambda^s v:\Lambda^s K_A^r\to\Lambda^s K_A^r
\]
的系数在 $A$ 中，且其模 $\pi_A$ 的归约
\[
v_s:\Lambda^s\kappa_A^r\to\Lambda^s\kappa_A^r
\]
不为零。若 $s\in\underline{r}^-\cup\underline{r}^+$，则 $v_s$ 的秩为 $1$，并分解为
\[
\begin{CD}
\Lambda^s\kappa_A^r @>{v_s}>> \Lambda^s\kappa_A^r\\
@VVV @AAA\\
\det(\kappa_A^r/F^s) @>{\sim}>> \det(E_s)
\end{CD}
\]
其中
\[
\begin{array}{rcl}
F^s & = & \text{$\kappa_A^r$ 中余维 $s$ 的子空间},\\
E_s & = & \text{$\kappa_A^r$ 中维数为 $s$ 的子空间}.
\end{array}
\]
此外，当 $s$ 在 $\underline{r}^-\cup\underline{r}^+$ 中变动时，
\[
\begin{array}{rcl}
(F^s) & = & \text{$\kappa_A^r$ 的下降滤过},\\
(E_s) & = & \text{$\kappa_A^r$ 的上升滤过}.
\end{array}
\]
对 $s\in\Bigl\{\begin{smallmatrix}\underline{r}^-\\ \underline{r}^+\end{smallmatrix}\Bigr\}$，把 $\begin{cases}s^+=\text{后继}\\ s^-=\text{前驱}\end{cases}$ 称为 $s$ 的后继与前驱。若 $s\in\underline{r}^-$，则 $v_{s+1}$ 分解为
\[
\begin{CD}
\Lambda^{s+1}\kappa_A^r @>{v_{s+1}}>> \Lambda^{s+1}\kappa_A^r\\
@VVV @AAA\\
\det(\kappa_A^r/F^s)\otimes(F^s/F^{s^+}) @>{\sim}>> \det(E_s)\otimes(E_{s^+}/E_s)
\end{CD}
\]

\origpage{22}由此得到诱导同构
\[
\overline{v}_s:F^s/F^{s^+}\xrightarrow{\sim}E_{s^+}/E_s.
\]

\par\medskip\noindent\textbf{基本注记.}\ 初等因子理论，因而前面的全部描述，对任意局部环 $A$ 都成立，只须假设：对每个 $s$，$u$ 的 $s$ 阶子式生成的理想是主理想。

\begin{definition*}
完整同态概形 $\Omega^r$ 指如下得到的概形：从 $M_r-\{0\}$ 出发，逐次爆炸“秩 $\le 1$ 的矩阵”“秩 $\le 2$ 的矩阵”、……、“秩 $\le r-1$ 的矩阵”所构成的闭子概形。
\end{definition*}

\begin{proposition*}
(i) $\Omega^r$ 以 $\mathrm{GL}_r$ 为稠密开集，并且带有 $\mathrm{GL}_r$ 的两个作用，一左一右。

(ii) $\Omega^r$ 光滑，其边界是正规交叉除子（DCN）。更确切地说，它带有一个 $\mathrm{GL}_r$-不变的光滑态射
\[
\Omega^r\to(\mathbb{A}^1/\mathbb{G}_m)^{r-1}.
\]

(ii)\footnote{译注：原文此处编号重复，两处均作 (ii)，原文如此。[译者按（上下文推断，无外部依据）：第二处或应作 (iii)。]} $\Omega^r$ 的边界分层 $\Omega^r_{\underline{r}}$ 以划分 $\underline{r}=(r_1+\cdots+r_k=r)$ 为下标。它们有模的描述。
\end{proposition*}

由 (i)，可以对概形 $S$ 上秩 $r$ 的两个丛 $\mathcal{F},\mathcal{E}$ 谈论“完整同态”
\[
\mathcal{F}\Rightarrow\mathcal{E}
\]
它由下列各项构成：

\begin{itemize}
\item 一族可逆层 $\mathcal{L}_s$，$1\le s<r$，各带 $S$ 上的整体截面 $\ell_s$，
\item 处处非零的同态
\[
v_s:\Lambda^s\mathcal{F}\otimes\bigotimes_{t<s}\mathcal{L}_t^{\otimes(s-t)}\to\Lambda^s\mathcal{E},\qquad 1\le s\le r,
\]
它们特别满足关系
\[
\Lambda^s v_1=\Bigl(\prod_{t<s}\ell_t^{(s-t)}\Bigr)\cdot v_s,\qquad 1\le s\le r.
\]
\end{itemize}

\origpage{23}这样的完整同态称为型 $\underline{r}=(r=r_1+\cdots+r_k)$ 的，如果：当 $s\in\underline{r}^-\cap\underline{r}^+$ 时 $\ell_s=0$，否则 $\ell_s$ 可逆。它由下列各项构成：

\begin{itemize}
\item 一族可逆层 $\mathcal{L}_s$，$s\in\underline{r}^-\cap\underline{r}^+$，
\item $\mathcal{F}$ 的一个下降滤过 $(\mathcal{F}^s)_{s\in\underline{r}^-\cup\underline{r}^+}$，由秩 $r-s$ 的极大子丛给出，
\item $\mathcal{E}$ 的一个上升滤过 $(\mathcal{E}_s)_{s\in\underline{r}^-\cup\underline{r}^+}$，由秩 $s$ 的极大子丛给出，
\item 同构
\[
\mathcal{F}^s/\mathcal{F}^{s^+}\otimes\Bigl(\bigotimes_{t\le s}\mathcal{L}_t\Bigr)\xrightarrow{\sim}\mathcal{E}_{s^+}/\mathcal{E}_s.
\]
\end{itemize}

\par\medskip
\noindent\textbf{迭代 chtouca 的叠}
\par\medskip

\begin{definition*}
设 $S$ 是 $\mathbb{F}_q$ 上的概形。$S$ 上秩 $r$ 的预迭代 chtouca 由下列各项构成：

\begin{itemize}
\item $X\times S$ 上秩 $r$ 的丛 $\mathcal{E}$，
\item $\mathcal{E}$ 的一个修正 $\mathcal{E}\hookrightarrow\mathcal{E}'\hookleftarrow\mathcal{E}''$（其中 $\mathcal{E}'/\mathcal{E}$ 与 $\mathcal{E}'/\mathcal{E}''$ 支撑于两个态射 $\infty,0:S\to X$ 的图像上，并在 $\mathcal{O}_S$ 上可逆），
\item $S$ 上的可逆层 $\mathcal{L}_1,\ldots,\mathcal{L}_{r-1}$，各带整体截面 $\ell_1,\ldots,\ell_{r-1}$，
\item 一个完整同态
\[
{}^\tau\mathcal{E}\Rightarrow\mathcal{E}''
\]
它在 $X\times S$ 上相伴的带截面可逆层依次为 $(\mathcal{L}_1,\ell_1)^{\otimes(q-1)},\ldots,(\mathcal{L}_{r-1},\ell_{r-1})^{\otimes(q-1)}$。
\end{itemize}
\end{definition*}

\origpage{24}设 $\mathcal{C}^r$ 为预迭代 chtouca 的（代数）叠。它带有态射
\[
\mathcal{C}^r\to(X\times X)\times(\mathbb{A}^1/\mathbb{G}_m)^{r-1}.
\]
记 $\mathcal{C}^r_{\underline{r}}$ 为它的各分层（即 $(\mathbb{A}^1/\mathbb{G}_m)^{r-1}$ 的点的原像）。设 $(\mathcal{E}\hookrightarrow\mathcal{E}'\hookleftarrow\mathcal{E}''\hookleftarrow{}^\tau\mathcal{E}\,;\,(\mathcal{L}_s,\ell_s)_{1\le s<r})$ 是型 $\underline{r}$ 的一个预迭代 chtouca。
完整同态 ${}^\tau\mathcal{E}\Rightarrow\mathcal{E}''$ 诱导两个由极大子丛给出的滤过：

\begin{itemize}
\item $(\mathcal{E}''_s)=$ $\mathcal{E}''$ 的上升滤过，
\item $(\overline{\mathcal{E}}_s)=$ ${}^\tau\mathcal{E}$ 的下降滤过。
\end{itemize}

考虑下列条件：

(i) 记 $\mathcal{E}'_s=\mathcal{E}''_s$（当 $s<r$）、$\mathcal{E}'_r=\mathcal{E}'$，则所有商 $\mathcal{E}'/\mathcal{E}'_s$ 都在 $\mathcal{O}_{X\times S}$ 上局部自由。

(ii) 若 $s\in\underline{r}^+$，则 $\mathcal{E}'_s\to\mathcal{E}'/\mathcal{E}$ 为满射，于是其核 $\mathcal{E}_s\hookrightarrow\mathcal{E}$ 在 $\mathcal{O}_{X\times S}$ 上局部自由（并约定 $\mathcal{E}_0=0$）。

(iii) 对每个 $s\in\underline{r}^+\cap\underline{r}^-$，有
\[
\overline{\mathcal{E}}_s\cap{}^\tau\mathcal{E}_s=0\quad\text{在 }{}^\tau\mathcal{E}\text{ 中}
\]
以及
\[
\overline{\mathcal{E}}_s+{}^\tau\mathcal{E}_{s^+}={}^\tau\mathcal{E}.
\]

这些条件定义 $\mathcal{C}^r_{\underline{r}}$ 的一个开集 $\mathrm{Cht}^r_{\underline{r}}$。它们的并是 $\mathcal{C}^r$ 的一个开集 $\overline{\mathrm{Cht}}{}^r$：这就是迭代 chtouca 的叠。

\begin{itemize}
\item 属于 $\mathrm{Cht}^r_{\underline{r}}$ 的迭代 chtouca 由下列各项构成：

  \begin{itemize}
  \item $X\times S$ 上的丛 $\mathcal{E}$ 以及一个修正 $\mathcal{E}\hookrightarrow\mathcal{E}'\hookleftarrow\mathcal{E}''$，
  \item 如上各滤过 $(\overline{\mathcal{E}}_s),(\mathcal{E}''_s),(\mathcal{E}'_s),(\mathcal{E}_s)$，
  \item 底概形 $S$ 上的可逆层 $(\mathcal{L}_s)_{s\in\underline{r}^-\cap\underline{r}^+}$，
  \item 同构
\[
\overline{\mathcal{E}}_{s^-}/\overline{\mathcal{E}}_s\otimes\Bigl(\bigotimes_{t<s}\mathcal{L}_t^{\otimes q}\Bigr)\xrightarrow{\sim}\mathcal{E}''_s/\mathcal{E}''_{s^-}\otimes\Bigl(\bigotimes_{t<s}\mathcal{L}_t\Bigr).
\]
  \end{itemize}
\end{itemize}

\par\medskip\noindent\origpage{25}\textbf{引理.}\ 对每个 $s\in\underline{r}^-\cap\underline{r}^+$，商
\[
{}^\tau\mathcal{E}/(\overline{\mathcal{E}}_s\oplus{}^\tau\mathcal{E}_s)
\]
支撑在一个态射
\[
x_s:S\to X\qquad(\text{秩 }s\text{ 处的“退化子”})
\]
的图像上，并在 $\mathcal{O}_S$ 上可逆。

\begin{proof}
只需在 $S$ 为几何点时证明。已知 ${}^\tau\mathcal{E}/(\overline{\mathcal{E}}_s\oplus{}^\tau\mathcal{E}_s)$ 是挠的。此外，
\begin{align*}
\deg{}^\tau\mathcal{E}/\overline{\mathcal{E}}_s&=\sum_{t\le s}\deg\overline{\mathcal{E}}_{t^-}/\overline{\mathcal{E}}_t=\sum_{t\le s}\deg\mathcal{E}''_t/\mathcal{E}''_{t^-}\\
&=\Bigl(\sum_{t\le s}\deg\mathcal{E}_t/\mathcal{E}_{t^-}\Bigr)+1=\deg\mathcal{E}_s+1=\deg{}^\tau\mathcal{E}_s+1.
\end{align*}
\end{proof}

在 $r_1=0^+$，即划分 $\underline{r}$ 的第一个元素处，令
\[
\mathcal{A}_{r_1}=\mathcal{E}_{r_1}\hookrightarrow\mathcal{E},
\]
\[
\mathcal{A}'_{r_1}=\mathcal{E}'_{r_1}=\mathcal{E}''_{r_1}\hookrightarrow\mathcal{E}'\hookleftarrow\mathcal{E}''.
\]
复合态射
\[
{}^\tau\mathcal{A}_{r_1}\hookrightarrow{}^\tau\mathcal{E}\to{}^\tau\mathcal{E}/\overline{\mathcal{E}}_{r_1}\xrightarrow{\sim}\mathcal{E}''_{r_1}=\mathcal{A}'_{r_1}
\]
可以嵌入一张图表
\[
\begin{pmatrix}
\mathcal{A}_{r_1} & \hookrightarrow & \mathcal{A}'_{r_1}\\
{}^\tau\mathcal{A}_{r_1} & &
\end{pmatrix}
\]
它是秩 $r_1$、极点为 $\infty$、零点为 $x_{r_1}$ 的 chtouca。

对 $s\in\underline{r}^+$，$s>r_1$，令
\[
\mathcal{A}_s=\mathcal{E}_s/\mathcal{E}_{s^-}\cong\mathcal{E}'_s/\mathcal{E}'_{s^-}
\]
\[
\mathcal{A}'_s=\overline{\mathcal{E}}_{s^-}\cap{}^\tau\mathcal{E}_s=\operatorname{Ker}\bigl[\overline{\mathcal{E}}_{s^-}\oplus{}^\tau\mathcal{E}_s\to{}^\tau\mathcal{E}\bigr].
\]

\origpage{26}两个复合态射
\[
\mathcal{A}'_s\hookrightarrow{}^\tau\mathcal{E}_s\to{}^\tau\mathcal{E}_s/{}^\tau\mathcal{E}_{s^-}={}^\tau\mathcal{A}_s
\]
与
\begin{align*}
\mathcal{A}'_s\otimes\bigotimes_{t<s}\mathcal{L}_t^q&\hookrightarrow\overline{\mathcal{E}}_{s^-}\otimes\bigotimes_{t<s}\mathcal{L}_t^q\to\overline{\mathcal{E}}_{s^-}/\overline{\mathcal{E}}_s\otimes\bigotimes_{t<s}\mathcal{L}_t^q\\
&\xrightarrow{\sim}\mathcal{E}''_s/\mathcal{E}''_{s^-}\otimes\bigotimes_{t<s}\mathcal{L}_t\hookrightarrow\mathcal{A}_s\otimes\bigotimes_{t<s}\mathcal{L}_t
\end{align*}
可以嵌入同一张图表
\[
\begin{pmatrix}
& & \mathcal{A}_s\otimes\bigotimes_{t<s}\mathcal{L}_t\\
\mathcal{A}'_s\otimes\bigotimes_{t<s}\mathcal{L}_t^q & \hookrightarrow & {}^\tau\bigl(\mathcal{A}_s\otimes\bigotimes_{t<s}\mathcal{L}_t\bigr)
\end{pmatrix}
\]
它定义一个“左”chtouca，秩为 $s-s^-$，极点为 $x_{s^-}$，零点为 $x_s$（其中 $x_r=0$）。

\begin{proposition*}
如此定义的态射
\[
\mathrm{Cht}^r_{\underline{r}}\to\mathrm{Cht}^{r_1}\times_X{}^{r_2}\mathrm{Cht}\times_X\cdots\times_X{}^{r_k}\mathrm{Cht}
\]
有限、满且根式。
\end{proposition*}

\section{紧化上的截断与水平结构}

已经定义了
\[
\overline{\mathrm{Cht}}{}^r=\coprod_{d\in\mathbb{Z}}\overline{\mathrm{Cht}}{}^{r,d}=\text{$X$ 上“迭代 chtouca”的叠}\,;
\]
它是局部有限型的代数叠（但不分离），并带有态射
\[
\overline{\mathrm{Cht}}{}^r\to X\times X\times(\mathbb{A}^1/\mathbb{G}_m)^{r-1}.
\]
$(\mathbb{A}^1/\mathbb{G}_m)^{r-1}$ 的点的原像（以划分 $\underline{r}=(r=r_1+\cdots+r_k)$ 为下标）是局部闭分层 $\mathrm{Cht}^r_{\underline{r}}$。开分层是
\[
\mathrm{Cht}^r_{\underline{0}}=\mathrm{Cht}^r
\]

\origpage{27}其余各分层则带有一个有限、满且根式的态射
\[
\mathrm{Cht}^r_{\underline{r}}\to\mathrm{Cht}^{r_1}\times_X{}^{r_2}\mathrm{Cht}\times_X\cdots\times_X{}^{r_k}\mathrm{Cht}.
\]

我们要在 $\overline{\mathrm{Cht}}{}^r$ 上定义截断。设 $\tilde{\mathcal{E}}=(\mathcal{E}\hookrightarrow\mathcal{E}'\hookleftarrow\mathcal{E}'',(\mathcal{L}_s,\ell_s)_{1\le s<r},{}^{\tau}\mathcal{E}\Rightarrow\mathcal{E}'')$ 是分层 $\mathrm{Cht}^r_{\underline{r}}$ 的一个点，而 $(\mathcal{E}_s),(\mathcal{E}'_s),(\mathcal{E}''_s),(\overline{\mathcal{E}}_s)$ 是由 $\mathcal{E},\mathcal{E}',\mathcal{E}'',{}^\tau\mathcal{E}$ 诱导的各滤过。

\begin{definition*}
(i) 设 $(\mathcal{F},\mathcal{F}')$ 是 $\mathcal{E}$ 与 $\mathcal{E}'$ 的极大子丛构成的一对。称它为这样的点 $\tilde{\mathcal{E}}$ 的“好”子对象，若满足

\begin{itemize}
\item $\mathcal{F}\hookrightarrow\mathcal{F}'$，且 $\operatorname{rg}\mathcal{F}=\operatorname{rg}\mathcal{F}'$，
\item 存在 $s\in\underline{r}^+$ 使\footnote{译注：下式第二个包含关系中的 $\mathcal{F}_s$ 原文如此。[译者按（上下文推断，无外部依据）：按上下文疑为 $\mathcal{E}'_s$。]}
\[
\mathcal{E}_{s^-}\varsubsetneq\mathcal{F}\subseteq\mathcal{E}_s,\qquad\mathcal{E}'_{s^-}\varsubsetneq\mathcal{F}'\subseteq\mathcal{F}_s
\]
\item 嵌入
\[
\overline{\mathcal{E}}_{s^-}/\overline{\mathcal{E}}_s\cong\mathcal{E}''_s/\mathcal{E}''_{s^-}\hookrightarrow\mathcal{E}'_s/\mathcal{E}'_{s^-}
\]
把 ${}^\tau\mathcal{F}\cap\overline{\mathcal{E}}_{s^-}/{}^\tau\mathcal{F}\cap\overline{\mathcal{E}}_s$ 送入 $\mathcal{F}'\cap\mathcal{E}'_s/\mathcal{F}'\cap\mathcal{E}'_{s^-}$。
\end{itemize}

(ii) 这样的“好”子对象称为

\begin{itemize}
\item 类型 I，若 $s=r_1=0^+$，或若 ${}^\tau\mathcal{F}+\overline{\mathcal{E}}_{s^-}={}^\tau\mathcal{E}$，
\item 类型 II，若 $s>r_1=0^+$，且 ${}^\tau\mathcal{F}+\overline{\mathcal{E}}_{s^-}\subsetneq{}^\tau\mathcal{E}$。
\end{itemize}
\end{definition*}

考虑一个凸截断多边形 $p:[0,r]\to\mathbb{R}_+$
\[
\begin{tikzpicture}[scale=0.9]
\draw[->] (-0.4,0) -- (7.2,0);
\draw[->] (0,-0.5) -- (0,2.6);
\draw (0,0) -- (1.6,1.35) -- (3.1,1.95) -- (4.6,1.55) -- (6.2,0);
\fill (1.6,1.35) circle (1.6pt);
\fill (4.6,1.55) circle (1.6pt);
\node[below left] at (0,0) {$0$};
\node[below] at (6.2,0) {$r$};
\end{tikzpicture}
\]

% ============================================================
%  \cite 使用清单
%  本文件对应原书 pp. 19--27（PDF 第 20--28 页），此段原文没有方括号文献
%  引用，因此没有用到任何 \cite。全书约定为“姓氏小写+两位年份”，例如
%  \cite{lafforgue97}、\cite{lafforgue98}、\cite{lafforgue02}、
%  \cite{drinfeld74}。
%
%  待定术语（术语表未收，此处暂定，待全书统一）
%  - strate：分层（术语表只收“strate de bord / 边界分层”）
%  - DCN (diviseur à croisements normaux)：正规交叉除子；相对者译作相对正规交叉除子
%  - polygone convexe de troncature：凸截断多边形（术语表收“截断多边形”）
%  - bon sous-objet：“好”子对象
%  - type I / type II：类型 I / 类型 II
%  - chtouca «à gauche»：“左”chtouca
%  - schéma des homomorphismes complets：完整同态概形
%  - section globale：整体截面
%  - description modulaire：模的描述
% ============================================================
